\chapter*{Zusammenfassung}\label{chap:summary}
\phantomsection
\addcontentsline{toc}{chapter}{Zusammenfassung}
\medskip
\noindent
Die vorliegende Projektarbeit befasst sich mit der prototypischen Entwicklung eines Slicers für die nichtplanare, robotergestützte additive Fertigung.
Konventionelle Slicer zerlegen Bauteile in rein horizontale Schichten.
Diese Methode bringt inhärente Nachteile mit sich, wie richtungsabhängige Bauteileigenschaften sowie die Notwendigkeit von Stützkonstruktionen für den Druck von Überhängen.
Um diese Limitationen zu überwinden, nutzt diese Studienarbeit die sechs Freiheitsgrade eines Industrieroboters, um die Werkzeugorientierung und Bahnplanung räumlich flexibel an die Bauteilgeometrie anzupassen.

Zur Lösung der softwareseitigen Herausforderungen wurde eine modulare Architektur entworfen, die aus zwei entkoppelten Hauptkomponenten besteht.
Das Modul \acs{medusa} übernimmt dabei die geometrische Bahnplanung, während das Modul \acs{kronos} als Middleware die Schnittstelle zum Industrieroboter bildet.

Im Rahmen der Arbeit wurden vier spezifische Slicing-Strategien konzipiert und in unterschiedlicher Tiefe untersucht:

\begin{enumerate}
    \item \textbf{Planarer Algorithmus} nutzt konventionelle horizontale Schnittebenen und dient im Projekt als funktionale Referenz und Validierungsbaseline.
    \item \textbf{Base-Algorithmus} ist ein vollständig implementierter Multiachsen-Ansatz.
    Er erkennt Überhänge geometrisch und spaltet das Bauteil dynamisch in einen \emph{Stamm} und in \emph{Zweige} auf.
    Diese Zweige werden mit eigenen, lokalen Wachstumsrichtungen gedruckt, was eine stützstrukturfreie Fertigung von Überhängen ermöglicht.
    \item \textbf{Sphären-Algorithmus} ist ein rein konzeptioneller Entwurf, der das Bauteil in konzentrische Sphärenschalen um einen zentralen Startpunkt zerlegt.
    Er wurde nicht implementiert, da an den Übergängen zwischen den Sphären mechanische Schwachstellen durch zu kleine Bindungsflächen zu erwarten sind.
    \item \textbf{Kristall-Algorithmus} ist ein komplexer, feldgetriebener Ansatz.
    Er basiert auf der Lösung der Laplace-Gleichung, um ein harmonisches Skalarfeld im Bauteilinneren zu erzeugen.
    Die Druckschichten folgen als gekrümmte Iso-Flächen der natürlichen Geometrie des Bauteils, wodurch die Werkzeugorientierung kontinuierlich variiert.
\end{enumerate}

Die Evaluation des Gesamtsystems erfolgte simulativ anhand von drei Testgeometrien (Würfel, Doppelpyramide, T-Stück) in der RViz-Umgebung.
Die Ergebnisse belegen die prinzipielle Machbarkeit der etablierten Datenkette vom \acs{CAD}-Modell bis zur simulierten 6-Achs-Roboterbewegung.
Mit dem Base-Algorithmus konnten 90\degree-Überhänge erfolgreich und stützfrei geplant werden.
Auch der Kristall-Algorithmus demonstriert erfolgreich die Generierung stützfreier, vollständig nicht-planarer Werkzeugpfade.

Als kritische Herausforderungen für die zukünftige Weiterentwicklung identifiziert die Arbeit die Grenzen des Inverse-Kinematik-Solvers sowie die Notwendigkeit,
in Zukunft die reale Extrusionsphysik und eine ebenenübergreifende Branch-Kollisionsprüfung zu integrieren.
Insgesamt liefert die Arbeit einen wertvollen und flexibel erweiterbaren „Proof of Concept“ zur Schließung der softwareseitigen Lücke in der robotergestützten additiven Fertigung.
\sig{Ben Stahl}